\documentclass[12pt]{article}
\usepackage[utf8]{inputenc}
\usepackage[T1]{fontenc}
\usepackage{hyperref}
\title{Kwerba Mamberamo Language Resources for K-12 Curriculum Development}
\author{Compiled by Claude Sonnet 5.5 (claude-sonnet-5-5) for Andr\'as Kornai}
\date{}
\begin{document}
\maketitle

\section{Introduction}
Kwerba Mamberamo (ISO 639-3 xwr; Glottocode kwer1263) is a Papuan language of Papua, Indonesia, in the Kwerba (Tor-Kwerba) family. Wikipedia and ASJP cite about 300 speakers (1993). Glottolog rates it not endangered, while one source calls it endangered. The searches found almost no description of the language itself.

\section{Primary Educational Resources}
No primers, textbooks or school programmes were found.

\section{Linguistic and Reference Materials}
\begin{itemize}
\item Wikipedia, Kwerba Mamberamo language (a stub): \url{https://en.wikipedia.org/wiki/Kwerba_Mamberamo_language}
\item Glottolog entry: \url{https://glottolog.org/resource/languoid/id/kwer1263}
\item ASJP wordlist: \url{https://asjp.clld.org/languages/KWERBA_MAMBERAMO}
\item Numerals data: \url{https://lingweb.eva.mpg.de/channumerals/Kwerba-Mamberamo.htm}
\item New Guinea World page on the Kwerba family. It mentions Silzer (1976), a grammar of Air Mati, a related language and not Kwerba Mamberamo: \url{https://newguineaworld.linguistik.uzh.ch/families/northwest-new-guinea/foja-range/west-foja-range/apauwar-river/kwerba}
\end{itemize}

\section{Religious and Community Resources}
None found.

\section{Digital and Audio Resources}
None found.

\section{Curriculum Development Context}
Almost nothing is documented for Kwerba Mamberamo. Any curriculum would first require fieldwork-based documentation and an orthography developed with the community. The related Kwerba-family literature may offer comparative help.
\end{document}